\documentclass[11pt, twocolumn]{article}
\usepackage[margin=0.65in, columnsep=0.25in, bottom=0.45in]{geometry}
\usepackage{newtxtext}
\usepackage{titlesec}
\usepackage{enumitem}
\usepackage{xurl}
\usepackage{hyperref}
\usepackage{xcolor}

% Tight compacting
\setlist{nosep, leftmargin=*} 
\linespread{0.93} 
\setlength{\parskip}{2.5pt} 
\renewcommand{\footnotesize}{\tiny}

% Tight section formatting with visual distinction
\titleformat{\section}{\large\bfseries\scshape}{\thesection}{1em}{}[\vspace{1pt}\titlerule]
\titlespacing*{\section}{0pt}{6pt}{3pt}

\hypersetup{
    colorlinks=true,
    linkcolor=black,
    urlcolor=blue,
    citecolor=black,
    breaklinks=true,
}

\pagestyle{plain} 

% Compact Bibliography
\makeatletter
\renewenvironment{thebibliography}[1]
     {\section*{References}%
      \tiny
      \setlength{\topsep}{0pt}
      \list{\@biblabel{\@arabic\c@enumiv}}%
           {\settowidth\labelwidth{\@biblabel{#1}}%
            \leftmargin\labelwidth
            \advance\leftmargin\labelsep
            \itemsep0pt\parsep0pt\parskip0pt
            \usecounter{enumiv}%
            \let\p@enumiv\@empty
            \renewcommand\theenumiv{\@arabic\c@enumiv}}%
      \sloppy
      \clubpenalty4000
      \@clubpenalty \clubpenalty
      \widowpenalty4000%
      \sfcode`\.\@m}
     {\def\@noitemerr
       {\@latex@warning{Empty `thebibliography' environment}}%
      \endlist}
\makeatother

\begin{document}

\twocolumn[
  \begin{@twocolumnfalse}
    \noindent
    \textbf{\LARGE AI Reading Exercise: Class 9} \hfill \textbf{Daniel Hardesty Lewis} \\
    \noindent
    \textit{AI and City Services} \hfill March 24, 2026 \\
    \textbf{PLAN A6613: AI and the Future of Cities} \hfill Spring 2026
    \vspace{3pt}
    \hrule
    \vspace{6pt}
  \end{@twocolumnfalse}
]

\section*{Tool \& Process}
I used \textbf{Claude Opus 4.6} with extended thinking via the Antigravity IDE.\footnote{Verbatim prompt log, downloaded sources, line-by-line annotations, a cross-source synthesis matrix, and a claim-by-claim verification audit are archived at \url{https://github.com/dhardestylewis/plan-a6613-ai-reading-class-3/tree/main/week9}.} I adopted the \textit{source-first} methodology developed in Week 7: three primary sources representing distinct stakeholder positions (municipal government, investigative journalism, advocacy reporting) were downloaded and annotated line by line before the LLM was asked to synthesize. A synthesis matrix then mapped every claim in this document to a specific annotated line. This proved essential because an unconstrained prompt produced a generically optimistic summary of municipal AI adoption that obscured the liability mechanisms central to this case. Specifically, the LLM's initial synthesis framed the chatbot's errors as isolated technical glitches rather than as a structural consequence of deploying a probabilistic text model in a domain where citizens expect deterministic, legally binding answers.

\section*{Key Findings}
\noindent New York City's MyCity Business chatbot illustrates how deploying a generative model as an official compliance advisor can reverse its intended benefit, replacing bureaucratic friction with legal friction.

\textbf{\textit{1. The system replaced human navigation of 2,000 pages of municipal code.}}
In October 2023, Mayor Eric Adams launched the MyCity chatbot as the centerpiece of the city's AI Action Plan, a framework described as ``the first of its kind for a major U.S. city'' (NYC Office of the Mayor, 2023). The chatbot builds on the first phase of the MyCity portal, launched in March 2023 to help families access child care (NYC Office of the Mayor, 2023), extending the same AI-driven approach to business compliance. Powered by Microsoft Azure AI, the bot was designed to let business owners ``access trusted information from more than 2,000 NYC Business web pages'' on ``compliance with codes and regulations, available business incentives, and best practices to avoid violations and fines'' (Lecher, 2024). The city states on the page hosting the bot that it will review questions to improve answers and delete user data within 30 days (Lecher, 2024).

\textbf{\textit{2. The system hallucinated structurally illegal advice.}}
Investigative testing by The Markup revealed the chatbot routinely generated advice contradicting settled law. It told business owners employers could take worker tips, violating Labor Law Section 196-d; that landlords could refuse Section 8 voucher holders, despite source-of-income discrimination being illegal since 2008; and that cashless stores were permitted, contradicting a 2020 city law (Lecher, 2024). These were confident inversions of statute, not ambiguous interpretations. The interface acknowledged this in fine print: it ``may occasionally produce incorrect, harmful or biased content'' (Lecher, 2024). Users had no mechanism to distinguish hallucinated advice from accurate guidance.

\textbf{\textit{3. The pilot-program defense shifted liability from city to citizen.}}
Mayor Adams defended keeping the chatbot online even as he acknowledged its answers were ``wrong in some areas'' (Offenhartz, 2024). Rosalind Black, Citywide Housing Director at Legal Services NYC, tested the bot herself after being alerted to The Markup's findings and discovered additional false information on housing, including that tenant lockouts were legal and that no restrictions existed on rent amounts (Lecher, 2024). She warned that ``if this chatbot is not being done in a way that is responsible and accurate, it should be taken down'' (Lecher, 2024). Julia Stoyanovich, computer science professor and director of the Center for Responsible AI at New York University, called the deployment ``reckless and irresponsible,'' arguing the city was ``rolling out software that is unproven without oversight'' (Offenhartz, 2024). The result is an unfunded mandate of verification: users must independently verify every output of a tool designed to eliminate the need for independent verification.

\section*{Verification and Reflection}
The AI correctly identified the MyCity chatbot and retrieved real sources rather than fabricating references. Its factual accuracy on the timeline, vendor, and specific hallucinations was confirmed against all three annotated sources. Its limitation, like the MyCity chatbot's own, was \textit{framing}: left unconstrained, it neutralized the conflict into a generic narrative of technical error without naming the statutes (196-d, the 2008 anti-discrimination law, the 2020 cashless ban) that make these hallucinations structurally illegal rather than merely inaccurate. Recovering this specificity required returning to the annotated sources. For a planner, this distinction matters: the failure is not that an AI tool occasionally errs, but that a municipal government deployed a probabilistic text generator as an authoritative legal advisor, transferring regulatory risk to citizens least equipped to detect the error.

\begin{thebibliography}{9}
\bibitem{ap2024} Offenhartz, J. (2024, April 3). NYC's AI chatbot was caught telling businesses to break the law. The city isn't taking it down. \textit{AP News}. \url{https://apnews.com/article/new-york-city-chatbot-misinformation-6ebc71db5b770b9969c906a7ee4fae21}
\bibitem{markup2024} Lecher, C. (2024, March 29). NYC's AI Chatbot Tells Businesses to Break the Law. \textit{The Markup}. \url{https://themarkup.org/artificial-intelligence/2024/03/29/nycs-ai-chatbot-tells-businesses-to-break-the-law}
\bibitem{nyc2023} NYC Office of the Mayor. (2023, Oct. 16). Mayor Adams Releases First-of-Its-Kind Plan For Responsible Artificial Intelligence Use In NYC Government. \url{https://www.nyc.gov/office-of-the-mayor/news/777-23/mayor-adams-releases-first-of-its-kind-plan-responsible-artificial-intelligence-use-nyc}
\end{thebibliography}

\end{document}
